\noindent \textit{\textcolor{red!60!black}{
One small comment---my understanding (though I could be wrong) is that repurchases in the open market were not allowed post 1929. I imagine the main reason driving the decline in shares was reorganizations. If that is the case, it is not obvious one would want to include this if the point of the analysis is to estimate whether firms were using payout policy to preemptively give the cash to shareholders. That said, the results in Table 4 show that the effects are driven by dividends, not repurchases, so this is mostly about guiding modern readers on what firm policies were in the 1930s.
}}

\bigskip

\noindent We are glad to clarify this for modern readers, and our data allow us to speak to it directly. Two facts stand out. First, treasury stock does appear on the Moody's balance sheets in our sample: between a fifth and a third of firms report a nonzero treasury-stock position during 1930--1934, and net additions to treasury stock occur in roughly one in eight firm-years. We cannot tell from the balance sheet how these shares were acquired---open-market purchases, private transactions, exchanges, or settlements---so this observation does not by itself settle the institutional question the referee raises. Whatever their source, the amounts are small: the median net addition is about 0.5\% of total assets. Second, our net repurchase measure is constructed from split-adjusted changes in shares outstanding, and its variation over 1930--1934 is dominated by the issuance side: the number of shares outstanding is entirely unchanged in two-thirds of firm-years, roughly two-thirds of the measure's variation comes from episodes of net issuance, and the share reductions that do occur are modest, with a median of about 1\% of market value. The referee's economic point therefore stands regardless of the institutional details: nothing in this period resembles the discretionary buyback programs of modern data, and net repurchases were not a meaningful payout vehicle.

\medskip
\noindent We agree that modern readers would benefit from this context, and we have added a footnote to Section 5 stating these facts: treasury-stock additions occur but are small, the number of shares outstanding is typically unchanged, and the variation in the net repurchase measure comes mostly from share issuance rather than retirement.

\medskip
\noindent This does not affect any of our substantive conclusions. As the referee notes, Table 5 (Table 4 in the previous draft) shows that the payout effect is driven entirely by cash dividends (column 2 is significant), while the net repurchase coefficient (column 3) is statistically insignificant in 1933 and 1934. The channel relevant to our debt overhang interpretation---shareholders extracting value through discretionary distributions in anticipation of losing their claim on the assets---operates through dividends, not through changes in shares outstanding.
